% Part: applied-modal-logics
% Chapter: epistemic-logic
% Section: properties-accessibility

\documentclass[../../../include/open-logic-section]{subfiles}

\begin{document}

\olfileid[id]{aml}{el}{acc}

\olsection{Relasi Aksesibilitas dan Prinsip Epistemik}

Berdasarkan apa yang telah kita ketahui tentang korespondensi kerangka dalam
logika modal normal, kita mungkin ingin melihat bentuk !!{formula}s penciri
ketika diberi penafsiran epistemik. Kita telah mengatakan bahwa logika
epistemik biasanya ditafsirkan dalam S5. Jadi, mari kita lihat bagaimana
berbagai prinsip epistemik direpresentasikan dan bagaimana prinsip-prinsip
itu berkorespondensi dengan berbagai syarat kerangka.

Ingat dari logika modal normal bahwa !!{formula}s modal yang berbeda mencirikan
sifat relasi aksesibilitas yang berbeda. Tabel berikut memilih beberapa di
antaranya yang berkorespondensi dengan prinsip epistemik tertentu.

\begin{table}[t]
    \begin{tabular}{| p{.48\textwidth} || p{.48\textwidth} |}
      \hline
      {\emph{Jika $R$ bersifat \dots}} & {\emph{maka \dots benar
      dalam~$\mModel{M}$:}} \\
      \hline \hline

      & $\Knows (p \lif q) \lif (\Knows p \lif \Knows q)$
      \hfill \newline{(Ketertutupan)} \\
      \hline
      \emph{refleksif}: $\forall w Rww$
      & $\Knows p \lif p$ \hfill \newline{(Veridikalitas)} \\
      \hline
      \emph{transitif}: \newline
      $\forall u \forall v \forall w ((Ruv \land Rvw) \lif Ruw)$ &
      $\Knows p \lif \Knows \Knows p$ \hfill
           \newline{(Introspeksi Positif)} \\
      \hline
      \emph{Euklidean}: \newline
      $\forall w \forall u \forall v ((Rwu \land Rwv) \lif Ruv)$ &
      $\lnot \Knows p \lif \Knows \lnot \Knows p$ \hfill
        \newline{(Introspeksi Negatif)} \\
      \hline
    \end{tabular}
    \caption{Empat prinsip epistemik.}
    \ollabel{tab:four}
  \end{table}

Veridikalitas, yang berkorespondensi dengan aksioma~$T$, sering dianggap
sebagai prinsip yang paling tidak kontroversial karena merepresentasikan
pernyataan bahwa jika !!a{formula} diketahui, formula itu harus benar.
Ketertutupan, serta Introspeksi Positif dan Negatif, jauh lebih
diperdebatkan.

Ketertutupan, yang berkorespondensi dengan aksioma~$K$, merepresentasikan
gagasan bahwa pengetahuan seorang agen tertutup terhadap implikasi. Dalam
beberapa kasus, hal ini mungkin tampak masuk akal. Sebagai contoh, saya
mungkin tahu bahwa jika saya berada di Victoria, saya berada di Pulau
Vancouver. Dengan mengesampingkan skenario skeptis yang ganjil, saya memang
tahu bahwa saya berada di Victoria, dan ini semestinya juga menunjukkan
bahwa saya tahu bahwa saya berada di Pulau Vancouver. Jadi, dalam kasus ini,
ketertutupan logis pengetahuan saya mungkin terasa relatif intuitif. Di sisi
lain, kita tidak selalu memikirkan semua konsekuensi pengetahuan kita,
sehingga hal ini dapat menghasilkan akibat yang kurang intuitif dalam kasus
lain.

Introspeksi Positif, yang terkadang dikenal sebagai prinsip KK, kadang
dirumuskan sebagai pernyataan bahwa jika saya mengetahui sesuatu, saya tahu
bahwa saya mengetahuinya. Prinsip ini merupakan padanan epistemik aksioma~4.
Sejalan dengan itu, introspeksi negatif dirumuskan sebagai pernyataan bahwa
jika saya \emph{tidak} mengetahui sesuatu, saya tahu bahwa saya tidak
mengetahuinya; prinsip ini merupakan padanan aksioma~5. Keduanya tampaknya
memiliki contoh penyangkal yang cukup lazim, tempat saya tidak yakin apakah
saya mengetahui sesuatu yang sebenarnya memang saya ketahui.


\end{document}
